On the CR5, I learned that command frequency is not the same as control quality. Under non-blocking execution, new inference consumed an intermediate in-motion state before the preceding action chunk had finished, causing back-and-forth corrections. Blocking, short-horizon ServoP made execution relatively serial and preserved state-action consistency at a lower apparent frequency. This was a problem of execution timing and semantics, not controller timing. It left me seeking deeper tools for connecting computer control, multivariable behavior, robot sensing, and implementation. NTU's MSc in Computer Control \& Automation matches that gap.

For the August 2027 intake, subject to course availability, EE6203 Computer Control Systems would strengthen the framework behind my digital-control decisions, while EE6225 Multivariable Control Systems Analysis \& Design would help me reason about coupled system behavior. EE6221 Robotics \& Intelligent Sensors would connect those analyses to sensed robot state. If approved, the optional three-AU, coursework-only EE6008 Collaborative Research \& Development Project could provide a team setting to investigate one bounded question: under what state-sampling and completion rules can short action chunks produce coherent physical trajectories? I would bring experience implementing synchronized observations, policy serving, and robot execution, and use this training to build control interfaces for dependable robotics and embodied-AI R\&D.
